% open-logic.tex
% compiles to produce complete text using open-logic-dev style

\documentclass[../include/open-logic-part]{subfiles}

\begin{document}

\clearpage

\begin{editorial}
يُحمِّل هذا الملف جميع المواد المضمَّنة في مشروع المنطق المفتوح.
وتدل الملاحظات التحريرية من هذا القبيل، إن ظهرت، على أن الملف جرى
تجميعه من غير أي اعتبار للكيفية التي ستُعرض بها هذه المواد. فإذا أمكنك
قراءة هذا، فالأرجح أنه \emph{لا يُنصح} باعتماد ملف PDF هذا في التدريس
أو الدراسة.

يوفّر مشروع المنطق المفتوح آليات كثيرة يمكن بها توليد نص أنسب للتدريس
أو للتعلّم الذاتي. فعلى سبيل المثال، يتعامل النص، في الإعداد الافتراضي،
مع جميع المؤثرات المنطقية بوصفها رموزًا أولية، ويعالج في البراهين جميع
الحالات لكل واحد من تلك المؤثرات. ومن الأفضل بكثير أن تُترك بعض هذه الحالات
تمارينَ. ومشروع المنطق المفتوح كذلك عمل قيد التطوير. وسعيًا إلى حفز
التعاون والتحسين، تُدرج المواد حتى إن كانت لا تزال مجرد مسودة، أو كانت
تفتقر إلى تمارين، وما إلى ذلك. وستُستبعد هذه الأقسام من أي ملف PDF
يُنتج لمقرر دراسي.

للعثور على ملفات PDF أنسب للتدريس والدراسة، راجع
\href{http://builds.openlogicproject.org/}{المقررات النموذجية
  المتاحة على موقع مشروع المنطق المفتوح}. ولإنشاء نسختك الخاصة، يمكنك
أن تبدأ من
\href{https://github.com/OpenLogicProject/OpenLogic/tree/master/courses/sample}{ملف
  التشغيل النموذجي}، أو أن تنظر في مصادر الكتب الدراسية المشتقة لترى
أمثلة أكثر تعقيدًا وتقدمًا.
\end{editorial}

\olimport[sets-functions-relations]{sets-functions-relations-complete}

\olimport[propositional-logic]{propositional-logic}

\olimport[first-order-logic]{first-order-logic}

\olimport[model-theory]{model-theory}

\olimport[computability]{computability}

\olimport[turing-machines]{turing-machines}

\olimport[incompleteness]{incompleteness}

\olimport[second-order-logic]{second-order-logic}

\olimport[lambda-calculus]{lambda-calculus}

\olimport[many-valued-logic]{many-valued-logic}

\olimport[normal-modal-logic]{normal-modal-logic}

\olimport[applied-modal-logic]{applied-modal-logic}

\olimport[intuitionistic-logic]{intuitionistic-logic}

\olimport[counterfactuals]{counterfactuals}

\olimport[set-theory]{set-theory}

\olimport[methods]{methods}

\olimport[history]{history}

\olimport[reference]{reference}

\end{document}
